
\section{Acknowledgements}
\label{sec:ack}

\dolma would not have been possible without the support of many individuals and institutions. 
The experimental components of this work were made possible through a partnership with AMD and CSC, enabling use of the LUMI supercomputer.
We thank Jonathan Frankle, Cody Blakeney, Matthew Leavitt and Daniel King and the rest of the MosaicML team for sharing findings from experiments on preliminary versions of our data.
We thank Vitaliy Chiley for messaging us on Twitter with a \href{https://twitter.com/vitaliychiley/status/1675594766799769600}{suggestion} for resolving a random number generator bug that was affecting our data shuffling.
We thank Erfan Al-Hossami, Shayne Longpre, and Gregory Yauney for sharing findings from their own large-scale pretraining data experiments.
We thank Ce Zhang and Maurice Weber of Together AI for thoughtful discussion on open datasets and data distribution format.
We thank Stella Biderman and Aviya Skowron for discussions around data licensing and data processing framework.
We thank our teammates at AI2 Nicole DeCario, Matt Latzke, Darrell Plessas, Kelsey MacMillan, Carissa Schoenick, Sam Skjonsberg, and Michael Schmitz for their help with the website, design, internal and external communications, budgeting, and other activities that supported smooth progress on this project.
Finally, we also express gratitude for the helpful discussions and feedback from our teammates at AI2 and close collaborators, including Prithviraj (Raj) Ammanabrolu, Maria Antoniak, Chris Callison-Burch, Peter Clark, Pradeep Dasigi, Nicole DeCario, Doug Downey, Ali Farhadi, Suchin Gururangan, Sydney Levine, Maarten Sap, Ludwig Schmidt, Will Smith, Yulia Tsvetkov, and Daniel S. Weld.

\section{Author Contributions}
\label{sec:contrib}

Dolma would not be possible without the help of our many teammates and collaborators. Weekly project meetings, messaging apps and documentation were accessible for anyone at AI2. Major decisions about Dolma were often made in these channels, with exception for certain topics (e.g., legal, funding). While many were involved in the Dolma effort (see Acknowledgements~\S\ref{sec:ack}), the authors of this paper were those who owned and delivered a critical piece of the puzzle. We detail their contributions below (authors in alphabetical order):

Contributors to \textbf{data acquisition and source-specific data processing} include Akshita Bhagia, Dirk Groeneveld, Rodney Kinney, Kyle Lo, Dustin Schwenk, and Luca Soldaini. Everyone contributed to literature review on available sources and best practices and decisions around sources to pursue. Akshita Bhagia, Rodney Kinney, Dustin Schwenk, and Luca Soldaini handled the bulk of data acquisition and processing and ablation experiments with 1B models for source-specific design decisions. Kyle Lo and Luca Soldaini handled discussions with legal to inform our choice of sources.

Contributors to \textbf{infrastructure and tooling} include Russell Authur, Dirk Groeneveld, Rodney Kinney, Kyle Lo, and Luca Soldaini. Rodney Kinney, Kyle Lo, and Luca Soldaini designed and implemented the shared toolkit used for processing our corpus at scale. Dirk Groeneveld wrote the Bloom filter for deduplication and decontamination. Russell Authur wrote a toolkit for acquisition and storage of Common Crawl data.

Contributors to \textbf{source-agnostic data processing} include Khyathi Chandu, Yanai Elazar, Rodney Kinney, Kyle Lo, Xinxi Lyu, Ian Magnusson, Aakanksha Naik, Abhilasha Ravichander, Zejiang Shen, and Luca Soldaini. Khyathi Chandu, and Aakanksha Naik developed the toxic text filter. Kyle Lo, and Xinxi Lyu helped evaluate it. Luca Soldaini developed the language filtering approach. Rodney Kinney, Zejiang Shen, and Luca Soldaini developed the ``quality'' filter. Yanai Elazar identified repeating $n$-gram sequences. Abhilasha Ravichander, Kyle Lo, and Luca Soldaini developed the PII filter. Jesse Dodge and Ian Magnusson developed the evaluation set decontamination approach. 

Contributors to \textbf{ablation experiments} include Iz Beltagy, Akshita Bhagia, Jesse Dodge, Dirk Groeneveld, Rodney Kinney, Kyle Lo, Ian Magnusson, Matthew Peters, Kyle Richardson, Dustin Schwenk, Luca Soldaini, Nishant Subramani, Oyvind Tafjord, and Pete Walsh. This work included designing and prioritizing experiments given compute constraints, implementing and running the 1B model experiments, and interpreting results. In particular, Oyvind Tafjord's work on the evaluation toolkit and Pete Walsh's work on the model implementation were critical.

Contributors to \textbf{posthoc experiments and analysis} on the final Dolma artifacts. Ben Bogin led the probing experiments on 1B model weights to assess impact of differing code mixtures with support from Kyle Lo and Niklas Muennighoff. Yanai Elazar ran the data analysis tool to summarize and document Dolma's composition. Valentin Hofmann led the tokenization fertility analysis with support from Kyle Lo. 
Ananya Harsh Jha and Ian Magnusson performed experiments training and evaluating baseline 1B models on other open datasets with support from Luca Soldaini.
Sachin Kumar and Jacob Morrison performed analysis of systematic issues in our choice of language identification and toxicity classifiers with support from Kyle Lo.
Niklas Muennighoff led analysis of correlation between different filters employed on Common Crawl data with support from Kyle Lo and Luca Soldaini.

Contributors to \textbf{licensing and release policy} include David Atkinson, Jesse Dodge, Jennifer Dumas, Nathan Lambert, Kyle Lo, Crystal Nam, and Luca Soldaini. 
David Atkinson, Jesse Dodge, Jennifer Dumas, and Crystal Nam led the bulk of this, including research into data licenses, risk-level determination for pretraining data, and defining the release policy. 
Kyle Lo and Luca Soldaini provided feedback throughout this process and handled technical details needed for the release.
Nathan Lambert provided feedback on release process and handled the actual release strategy, particularly around external communication.

All of the contributors above helped with \textbf{documentation and writing} of their respective components. In particular, Li Lucy provided an extensive literature review of language models, open corpora and pretraining corpus creation practices. Emma Strubell gave valuable feedback on our manuscript.
Nathan Lambert helped with feedback on the blog post and other forms of external-facing communication about Dolma.

Hannaneh Hajishirzi, Noah Smith, and Luke Zettlemoyer \textbf{advised} on the project, including broad strategy, writing, recruiting and providing resources. As OLMo project leads, Iz Beltagy, Jesse Dodge, and Dirk Groeneveld helped with \textbf{visibility and coordination} with other critical OLMo project workstreams. Notably, we credit Noah Smith for coming up with the name Dolma.

Finally, Kyle Lo and Luca Soldaini \textbf{led} the overall Dolma project and were involved in all aspects, including project management, planning and design, discussions with legal and ethics committees, data and compute partnerships, infrastructure, tooling, implementation, experiments, writing/documentation, etc.
